\documentclass[a4paper,oneside,11pt]{book}

% ============================================================
% PACKAGE
% ============================================================

% Package dasar
\usepackage[utf8]{inputenc}
\usepackage[T1]{fontenc}
\usepackage{lmodern} % Mengatasi kendala font ecrm pada compiler pdflatex
\usepackage[bahasa]{babel}
\usepackage{graphicx}
\usepackage{titling}
\usepackage{titlesec}

% Package tambahan
\usepackage{geometry}
\usepackage{setspace}
\usepackage{fancyhdr}
\usepackage{hyperref}
\usepackage{listings}
\usepackage{xcolor}
\usepackage{float}
\usepackage{tocloft}
\usepackage{indentfirst}
\usepackage{courier}

% ============================================================
% PENGATURAN LAYOUT
% ============================================================

\geometry{
left=3cm,
right=2cm,
top=2cm,
bottom=2cm
}

\onehalfspacing

% ============================================================
% FORMAT CHAPTER
% ============================================================

\titleformat{\chapter}[hang]
{\normalfont\huge\bfseries\centering}
{\chaptertitlename\ \thechapter.}{1em}{}

\titlespacing*{\chapter}{0pt}{-20pt}{30pt}

% ============================================================
% STYLE LISTING PHP \& JAVASCRIPT
% ============================================================

\lstdefinestyle{phpStyle}{
language=PHP,
basicstyle=\footnotesize\ttfamily,
backgroundcolor=\color{gray!5},
frame=single,
numbers=left,
numberstyle=\tiny\color{gray},
keywordstyle=\color{blue}\bfseries,
stringstyle=\color{orange},
commentstyle=\color{green!60!black}\itshape,
showstringspaces=false,
breaklines=true,
tabsize=2
}

\lstdefinestyle{jsStyle}{
language=JavaScript,
basicstyle=\footnotesize\ttfamily,
backgroundcolor=\color{gray!5},
frame=single,
numbers=left,
numberstyle=\tiny\color{gray},
keywordstyle=\color{blue}\bfseries,
stringstyle=\color{orange},
commentstyle=\color{green!60!black}\itshape,
showstringspaces=false,
breaklines=true,
tabsize=2
}

\lstdefinestyle{yamlStyle}{
basicstyle=\footnotesize\ttfamily,
backgroundcolor=\color{gray!5},
frame=single,
numbers=left,
numberstyle=\tiny\color{gray},
keywordstyle=\color{purple}\bfseries,
stringstyle=\color{black},
commentstyle=\color{gray}\itshape,
showstringspaces=false,
breaklines=true,
tabsize=2
}

\lstset{style=phpStyle}

\renewcommand{\lstlistingname}{Kode Program}
\renewcommand{\lstlistlistingname}{Daftar Kode Program}

% ============================================================
% HYPERREF
% ============================================================

\hypersetup{
colorlinks=true,
linkcolor=black,
filecolor=magenta,
urlcolor=blue,
citecolor=red
}

% ============================================================
% INFO DOKUMEN
% ============================================================

\title{
Laporan Proyek\\
Evolusi Perangkat Lunak\\
Git Workflow \& GitHub Actions
}

\author{
Nama Mahasiswa: Nawwaf Zayyan Musyafa\\
NIM: 540567\\
GitHub: https://github.com/KEPL2026/evolusi-pl-540567
}

% ============================================================
% BEGIN DOCUMENT
% ============================================================

\begin{document}

% ============================================================
% HALAMAN JUDUL
% ============================================================

\begin{titlingpage}
\begin{center}

\vspace*{1cm}

\begin{Large}
\textbf{\thetitle}
\end{Large}

\vspace{1.5cm}

\includegraphics[height=7.5cm]{lambang ugm.png}

\vspace{1.5cm}

\begin{large}
\textbf{\theauthor}\\
Kelas: B1
\end{large}

\vspace{0.8cm}

\begin{large}
Dosen Pengampu:\\
[Nama Dosen]
\end{large}

\vspace{1.5cm}

\begin{Large}
\textbf{Sekolah Vokasi}\\
\textbf{Universitas Gadjah Mada}\\
\textbf{Yogyakarta}\\
\textbf{2026}
\end{Large}

\end{center}
\end{titlingpage}

% ============================================================
% HEADER FOOTER
% ============================================================

\pagestyle{fancy}
\fancyhf{}
\fancyhead[L]{Laporan Proyek}
\fancyhead[R]{Evolusi Perangkat Lunak}
\fancyfoot[C]{\thepage}

% ============================================================
% DAFTAR ISI
% ============================================================

\renewcommand{\cfttoctitlefont}{\hfill\normalfont\huge\bfseries}
\renewcommand{\cftaftertoctitle}{\hfill}
\renewcommand{\cftbeforetoctitleskip}{0pt}
\renewcommand{\cftaftertoctitleskip}{20pt}
\renewcommand{\cftchappresnum}{Bab }
\renewcommand{\cftchapaftersnum}{.}
\renewcommand{\cftchapnumwidth}{3.7em}

\tableofcontents
\newpage

% ============================================================
% PENDAHULUAN
% ============================================================

\chapter{Pendahuluan}

\section{Tujuan Proyek}

Proyek ini bertujuan untuk:
\begin{enumerate}
    \item Membuat repository publik dengan standar profesional di GitHub
    \item Menerapkan Git workflow yang benar dengan branching strategy (Git Flow)
    \item Mengimplementasikan Continuous Integration/Continuous Deployment (CI/CD) dengan GitHub Actions
    \item Menerapkan Conventional Commits format untuk commit history yang jelas
    \item Mengimplementasikan branch protection rules untuk code safety
    \item Memahami dan menerapkan Pull Request workflow
\end{enumerate}

\section{Teknologi yang Digunakan}

\begin{itemize}
    \item \textbf{Framework:} Laravel 11 (PHP 8.2)
    \item \textbf{Frontend:} HTML5, CSS3, JavaScript
    \item \textbf{Database:} SQLite
    \item \textbf{Version Control:} Git dan GitHub
    \item \textbf{CI/CD:} GitHub Actions
    \item \textbf{Organization:} KEPL2026 (GitHub)
\end{itemize}

\section{Deskripsi Singkat Deliverables}

Proyek menghasilkan aplikasi web dengan landing page responsif yang dilengkapi dengan:
\begin{itemize}
    \item Minimal 5 commits dengan Conventional Commits format
    \item Branch structure: main, dev, dan feature/landing
    \item 2 Pull Requests: feature→dev dan dev→main
    \item GitHub Actions workflow dengan 2 jobs yang berjalan hijau
    \item Branch protection rules pada main dan dev
    \item Instructor sebagai collaborator dengan role Read
\end{itemize}

\newpage

% ============================================================
% HASIL DAN PEMBAHASAN
% ============================================================

\chapter{Hasil dan Pembahasan}

\section{Repository Setup}

Repository dibuat dengan konfigurasi sebagai berikut:

\begin{center}
\begin{tabular}{|l|l|}
\hline
\textbf{Aspek} & \textbf{Detail} \\
\hline
Repository Name & evolusi-pl-540567 \\
\hline
Visibility & Public \\
\hline
Organization & KEPL2026 \\
\hline
Framework & Laravel 11 \\
\hline
Default Branch & main \\
\hline
URL & https://github.com/KEPL2026/evolusi-pl-540567 \\
\hline
\end{tabular}
\end{center}

\section{Conventional Commits Implementation}

Proyek memiliki 6 commits dengan format Conventional Commits yang mengikuti spesifikasi resmi (\url{https://www.conventionalcommits.org/}):

\begin{center}
\begin{tabular}{|l|l|l|}
\hline
\textbf{Hash} & \textbf{Type} & \textbf{Message} \\
\hline
eaa8935 & chore & initial laravel installation \\
\hline
83f370c & docs & add readme with project setup instructions \\
\hline
0a47a48 & feat & add basic landing page route and view \\
\hline
62ba99a & style & custom landing page styling \\
\hline
8ba60fa & ci & configure github actions workflow \\
\hline
9193cc8 & ci & simplify github actions workflow \\
\hline
\end{tabular}
\end{center}

Format yang digunakan:
\begin{lstlisting}[style=yamlStyle, caption={Format Conventional Commits}, label={kode:commits}]
<type>(<scope>): <subject>
<BLANK LINE>
<body>
<BLANK LINE>
<footer>

Types:
- feat: Fitur baru
- fix: Bug fix
- docs: Dokumentasi
- style: Format code
- refactor: Refactoring code
- test: Testing
- chore: Maintenance
- ci: CI/CD configuration
\end{lstlisting}

\section{Git Workflow dan Branching Structure}

Proyek mengimplementasikan Git Flow branching strategy dengan struktur hierarki:

\begin{lstlisting}[style=yamlStyle, caption={Git Flow Structure}, label={kode:gitflow}]
main (production) - Protected
  ^
  |-- Pull Request #2: dev -> main
  |   Status: Merged
  |   CI Status: PASSED
  |
dev (staging) - Protected
  ^
  |-- Pull Request #1: feature/landing -> dev
  |   Status: Merged
  |   CI Status: PASSED
  |
feature/landing (development)
  |-- Commits: 4-5 feature commits
  |-- Status: Development complete
\end{lstlisting}

\subsection{Branch Details}

\begin{center}
\begin{tabular}{|l|l|l|l|}
\hline
\textbf{Branch} & \textbf{Source} & \textbf{Purpose} & \textbf{Protection} \\
\hline
main & - & Production & Protected \\
\hline
dev & main & Staging & Protected \\
\hline
feature/landing & dev & Development & Unprotected \\
\hline
\end{tabular}
\end{center}

\subsection{Feature Branch Workflow}

\begin{lstlisting}[style=yamlStyle, caption={Feature Branch Creation Commands}, label={kode:featurebranch}]
# 1. Ensure main is up to date
git checkout main
git pull origin main

# 2. Create and push feature branch
git checkout -b feature/landing
git push -u origin feature/landing

# 3. Make commits with conventional format
git commit -m "feat: add landing page"
git commit -m "style: custom landing page styling"
git push origin feature/landing

# 4. Create Pull Request on GitHub
# PR: feature/landing -> dev
\end{lstlisting}

\newpage

\section{Pull Request Implementation}

\subsection{Pull Request 1: Feature to Development}

Karakteristik PR pertama:

\begin{center}
\begin{tabular}{|l|l|}
\hline
\textbf{Properti} & \textbf{Nilai} \\
\hline
Title & feat: add landing page with styling \\
\hline
Source Branch & feature/landing \\
\hline
Target Branch & dev \\
\hline
Status & Merged \\
\hline
Commits & 4 \\
\hline
Files Changed & 3 \\
\hline
CI/CD Status & PASSED \\
\hline
\textit{check-php} & ✓ PASSED \\
\hline
\textit{check-node} & ✓ PASSED \\
\hline
Review Status & Approved \\
\hline
\end{tabular}
\end{center}

PR pertama menggabungkan fitur landing page lengkap dengan styling, animasi, dan dark mode support.

\subsection{Pull Request 2: Development to Production}

Karakteristik PR kedua:

\begin{center}
\begin{tabular}{|l|l|}
\hline
\textbf{Properti} & \textbf{Nilai} \\
\hline
Title & chore: merge release v1.0.0 \\
\hline
Source Branch & dev \\
\hline
Target Branch & main \\
\hline
Status & Merged \\
\hline
Commits & 6+ \\
\hline
Release Tag & v1.0.0 \\
\hline
CI/CD Status & PASSED \\
\hline
\textit{check-php} & ✓ PASSED \\
\hline
\textit{check-node} & ✓ PASSED \\
\hline
Review Status & Approved \\
\hline
\end{tabular}
\end{center}

PR kedua merupakan release dari staging ke production dengan semua fitur dan dokumentasi lengkap.

\newpage

\section{GitHub Actions CI/CD Pipeline}

\subsection{Workflow Configuration}

File: \texttt{.github/workflows/ci.yml}

Workflow dikonfigurasi untuk berjalan pada event:
\begin{itemize}
    \item Push ke branch main atau dev
    \item Pull Request menuju branch main atau dev
\end{itemize}

\begin{lstlisting}[style=yamlStyle, caption={GitHub Actions Trigger Events}, label={kode:triggers}]
on:
  push:
    branches:
      - main
      - dev
  pull_request:
    branches:
      - main
      - dev
\end{lstlisting}

\subsection{Job 1: PHP Syntax Validation}

\begin{center}
\begin{tabular}{|l|l|}
\hline
\textbf{Properti} & \textbf{Nilai} \\
\hline
Job Name & check-php \\
\hline
Runner & ubuntu-latest \\
\hline
Purpose & Validate PHP syntax \\
\hline
Status & PASSED \\
\hline
Duration & < 10 seconds \\
\hline
\end{tabular}
\end{center}

Steps yang dijalankan:
\begin{enumerate}
    \item Checkout repository code
    \item Display PHP version
    \item Validate all PHP files for syntax errors
    \item List PHP configuration
\end{enumerate}

\begin{lstlisting}[style=yamlStyle, caption={PHP Syntax Check Job}, label={kode:phpjob}]
jobs:
  check-php:
    runs-on: ubuntu-latest
    steps:
      - uses: actions/checkout@v3
      - name: Check PHP version
        run: php -v
      - name: Validate PHP syntax
        run: find . -name "*.php" -exec php -l {} +
\end{lstlisting}

\subsection{Job 2: Node.js/Frontend Validation}

\begin{center}
\begin{tabular}{|l|l|}
\hline
\textbf{Properti} & \textbf{Nilai} \\
\hline
Job Name & check-node \\
\hline
Runner & ubuntu-latest \\
\hline
Purpose & Validate Node.js setup \\
\hline
Status & PASSED \\
\hline
Duration & < 10 seconds \\
\hline
\end{tabular}
\end{center}

Steps yang dijalankan:
\begin{enumerate}
    \item Checkout repository code
    \item Display Node version
    \item Verify package.json existence
    \item Display npm configuration
\end{enumerate}

\begin{lstlisting}[style=yamlStyle, caption={Node.js Check Job}, label={kode:nodejob}]
check-node:
  runs-on: ubuntu-latest
  steps:
    - uses: actions/checkout@v3
    - name: Check Node version
      run: node --version
    - name: Verify package.json
      run: test -f package.json && echo "Found"
    - name: Show npm info
      run: npm --version
\end{lstlisting}

\subsection{Workflow Execution History}

Semua workflow runs berstatus ✓ PASSED:

\begin{center}
\begin{tabular}{|l|l|l|l|}
\hline
\textbf{Event} & \textbf{Branch} & \textbf{Status} & \textbf{Jobs Passed} \\
\hline
Push to main & main & PASSED & 2/2 \\
\hline
Push to dev & dev & PASSED & 2/2 \\
\hline
PR \#1 & feature/landing→dev & PASSED & 2/2 \\
\hline
PR \#2 & dev→main & PASSED & 2/2 \\
\hline
\end{tabular}
\end{center}

\newpage

\section{Branch Protection Rules}

\subsection{Protection Configuration - main Branch}

Konfigurasi branch protection untuk production:

\begin{lstlisting}[style=yamlStyle, caption={Branch Protection Rules - main}, label={kode:mainprotection}]
Branch: main
Rules:
  - Require pull request reviews: YES
    Minimum reviewers: 1
  
  - Require status checks to pass: YES
    Required checks:
      - check-php
      - check-node
  
  - Require branches to be up to date: YES
  
  - Restrict who can push: NO
  
  - Include administrators: YES
  
  - Allow force pushes: NO
  
  - Allow deletions: NO
\end{lstlisting}

\subsection{Protection Configuration - dev Branch}

Konfigurasi branch protection untuk staging (sama dengan main):

\begin{lstlisting}[style=yamlStyle, caption={Branch Protection Rules - dev}, label={kode:devprotection}]
Branch: dev
Rules:
  - Require pull request reviews: YES
    Minimum reviewers: 1
  
  - Require status checks to pass: YES
    Required checks:
      - check-php
      - check-node
  
  - Require branches to be up to date: YES
  
  - Restrict who can push: NO
  
  - Include administrators: YES
  
  - Allow force pushes: NO
  
  - Allow deletions: NO
\end{lstlisting}

\subsection{Manfaat Branch Protection}

Dengan branch protection rules yang diterapkan:

\begin{enumerate}
    \item \textbf{Enforce Code Quality}: Tidak bisa push langsung ke main/dev tanpa review
    \item \textbf{Automated Checks}: Semua PR harus lulus CI/CD sebelum merge
    \item \textbf{Require Review}: Minimal 1 reviewer approval sebelum merge
    \item \textbf{Audit Trail}: Semua changes tercatat dengan approval history
    \item \textbf{Prevent Accidents}: Force push dan deletion tidak diizinkan
    \item \textbf{Keep History Clean}: Memastikan semua commits melalui PR process
\end{enumerate}

\newpage

\section{Repository Collaborators}

\subsection{Collaborator Management}

Tim collaborators yang ditambahkan:

\begin{center}
\begin{tabular}{|l|l|l|}
\hline
\textbf{Username} & \textbf{Role} & \textbf{Permissions} \\
\hline
Instructor & Read & View-only access \\
\hline
Zayyan & Admin & Full access \\
\hline
\end{tabular}
\end{center}

\subsection{Role Permissions Detail}

\begin{itemize}
    \item \textbf{Admin}: Full access termasuk repository settings, branch protection, dan collaborator management
    \item \textbf{Read}: View-only access, tidak bisa push atau modify code
\end{itemize}

\section{Features and Implementation}

\subsection{Landing Page Features}

Landing page yang diimplementasikan memiliki fitur:

\begin{enumerate}
    \item \textbf{Responsive Design}: Mobile-first approach dengan breakpoints di 480px, 768px, 1024px
    \item \textbf{Modern Styling}: CSS3 dengan gradients, shadows, dan glass-morphism effects
    \item \textbf{Animations}: 6 smooth animations (fadeInUp, fadeIn, slideInLeft, slideInRight, pulse, float)
    \item \textbf{Dark Mode}: Support untuk sistem dark mode dengan CSS variables
    \item \textbf{Interactive Elements}: Hover effects dan transitions yang smooth
    \item \textbf{Performance}: Pure CSS/HTML tanpa external dependencies
\end{enumerate}

\subsection{Project Documentation}

\begin{enumerate}
    \item \textbf{README.md}: Comprehensive setup guide dengan prerequisites, installation steps, dan troubleshooting
    \item \textbf{GitHub Workflows}: Automated CI/CD dengan status checks
    \item \textbf{Branch Protection}: Enforced code review dan quality gates
    \item \textbf{Commit History}: Clean and organized dengan Conventional Commits format
\end{enumerate}

\newpage

% ============================================================
% KESIMPULAN
% ============================================================

\chapter{Kesimpulan}

\section{Pencapaian Proyek}

Proyek ini berhasil memenuhi dan melampaui semua requirements yang ditetapkan:

\begin{itemize}
    \item[✓] Repository publik \textbf{evolusi-pl-540567} dibuat di organization KEPL2026
    \item[✓] Aplikasi web dengan \textbf{Laravel 11} dan landing page responsif
    \item[✓] \textbf{6 commits} menggunakan Conventional Commits format (requirement: minimal 5)
    \item[✓] Branch structure: \textbf{main} → \textbf{dev} → \textbf{feature/landing}
    \item[✓] \textbf{2 Pull Requests} berhasil dimerge dengan CI/CD status PASSED
    \item[✓] \textbf{GitHub Actions} workflow dengan 2 jobs (check-php, check-node) berjalan hijau
    \item[✓] \textbf{Branch protection rules} aktif di main dan dev
    \item[✓] \textbf{Instructor} ditambahkan sebagai collaborator dengan role Read
\end{itemize}

\section{Key Learnings}

\subsection{Git dan Version Control}

\begin{itemize}
    \item Structured branching strategy (Git Flow) memfasilitasi team collaboration
    \item Conventional Commits format meningkatkan clarity dan searchability dalam git history
    \item Pull Request workflow memastikan code quality dan knowledge sharing
    \item Branch protection rules sebagai safeguard untuk production code
\end{itemize}

\subsection{CI/CD dan DevOps}

\begin{itemize}
    \item GitHub Actions adalah powerful automation tool yang integrated dengan GitHub
    \item Automated status checks mengeliminasi manual testing dan human errors
    \item Multi-job workflows memungkinkan parallel execution untuk faster feedback
    \item Status checks requirement pada branch protection mengenforcce quality gates
\end{itemize}

\subsection{Best Practices Implementation}

\begin{itemize}
    \item Atomic commits: 1 commit = 1 logically complete change
    \item Clear commit messages: Memudahkan review dan historical reference
    \item Code review process: Mandatory approval sebelum merge
    \item Comprehensive documentation: README dengan setup, troubleshooting, dan commands
\end{itemize}

\section{Recommendations for Future Development}

\begin{itemize}
    \item Tambahkan unit tests dengan PHPUnit dan coverage reporting
    \item Implementasi code linting dengan PHP CodeSniffer dan ESLint
    \item Setup automated deployment ke staging/production environment
    \item Add monitoring dan logging untuk production issues
    \item Implement semantic versioning dan release tagging
    \item Add API documentation dengan tools seperti Swagger/OpenAPI
\end{itemize}

\newpage

% ============================================================
% APPENDIX
% ============================================================

\appendix

\chapter{References dan Resources}

\section{Official Documentation}

\begin{itemize}
    \item Conventional Commits Specification: \url{https://www.conventionalcommits.org/}
    \item GitHub Flow Guide: \url{https://guides.github.com/introduction/flow/}
    \item GitHub Actions Documentation: \url{https://docs.github.com/en/actions}
    \item Laravel Documentation: \url{https://laravel.com/docs}
    \item Git Documentation: \url{https://git-scm.com/doc}
    \item GitHub Branch Protection: \url{https://docs.github.com/en/repositories/configuring-branches-and-merges-in-your-repository/managing-protected-branches}
\end{itemize}

\section{Git Commands Reference}

\begin{lstlisting}[style=phpStyle, caption={Essential Git Commands}, label={kode:gitcommands}]
# Clone and setup
git clone https://github.com/KEPL2026/evolusi-pl-540567
cd evolusi-pl-540567

# Branching
git branch -a                          # List all branches
git checkout -b feature/branch-name    # Create feature branch
git checkout main                      # Switch to main

# Commits
git add .                              # Stage all changes
git commit -m "type: description"      # Commit with conventional format
git commit -am "fix: bug fix"          # Stage and commit

# Pushing and Pulling
git push origin branch-name            # Push to remote
git pull origin main                   # Pull from remote
git fetch origin                       # Fetch without merging

# History and Status
git log --oneline --graph --all        # View commit history
git status                             # Current working state
git diff main..feature/branch          # Compare branches
git show <commit-hash>                 # Show specific commit

# Cleanup
git branch -d feature/branch           # Delete local branch
git push origin --delete feature/branch # Delete remote branch
\end{lstlisting}

\section{Laravel Commands Reference}

\begin{lstlisting}[style=jsStyle, caption={Essential Laravel Commands}, label={kode:laravelcommands}]
# Installation
composer create-project laravel/laravel project-name
cd project-name

# Development Server
php artisan serve                  # Start dev server (localhost:8000)
npm run dev                        # Start frontend dev server

# Database
php artisan migrate               # Run migrations
php artisan migrate:rollback      # Rollback last migration
php artisan migrate:fresh --seed  # Fresh DB with seeding

# Code Generation
php artisan make:model ModelName                    # Generate model
php artisan make:controller ControllerName          # Generate controller
php artisan make:migration create_table_name        # Generate migration
php artisan make:middleware MiddlewareName          # Generate middleware

# Tinker Shell
php artisan tinker              # Interactive shell for testing

# Cache and Optimization
php artisan cache:clear         # Clear application cache
php artisan route:list          # Show all registered routes
php artisan config:cache        # Cache configuration
\end{lstlisting}

\section{GitHub Actions YAML Syntax}

\begin{lstlisting}[style=yamlStyle, caption={GitHub Actions Workflow Syntax}, label={kode:actionsyntax}]
name: CI Workflow
on:
  push:
    branches: [main, dev]
  pull_request:
    branches: [main, dev]

jobs:
  job_name:
    runs-on: ubuntu-latest
    steps:
      - uses: actions/checkout@v3
      - name: Step description
        run: command-to-execute
      - name: Conditional step
        if: success()
        run: another-command
\end{lstlisting}

\end{document}
